%%%%%%%%%%%%%%%%%%%%%%%%%%%%%%%%%%%%%%%%%%%%%%%%%%%%%%%%%%%%%%%%%%%%%%%
% BAB 1
%%%%%%%%%%%%%%%%%%%%%%%%%%%%%%%%%%%%%%%%%%%%%%%%%%%%%%%%%%%%%%%%%%%%%%%
\mychapter{1}{BAB 1 PENDAHULUAN}

\section{Latar Belakang}

Kabupaten Sukabumi merupakan salah satu wilayah dengan luas daratan terbesar di Jawa Barat \parencite{bps-luas-2025} dengan sektor Pertanian, Kehutanan, dan Perikanan sebagai kontributor terbesar terhadap perekonomian daerah dengan kontribusi sebesar 24,27\% dan secara konsisten menjadi sektor dengan kontribusi tertinggi selama periode 2021–2025 \parencite{bps-sukabumi-2026}. Aktivitas distribusi komoditas pangan antarkecamatan menjadi salah satu faktor penting dalam menjaga ketersediaan pasokan dan stabilitas harga bahan pokok. Dengan jumlah penduduk mencapai lebih dari 2,6 juta jiwa dan 47 kecamatan yang tersebar mulai dari wilayah pegunungan hingga pesisir selatan \parencite{bps-penduduk-2025}, stabilitas harga bahan pokok menjadi tantangan utama dalam menjaga kesejahteraan masyarakat.

Komoditas pangan strategis termasuk ke dalam kelompok volatile food, yaitu komoditas yang memiliki karakteristik harga cenderung berfluktuasi. Kelompok komoditas tersebut meliputi beras, daging ayam, daging sapi, telur ayam, bawang merah, bawang putih, cabai merah, cabai rawit, minyak goreng, dan gula pasir \parencite{Helbawanti_Saputro_Ulfa_2021}. Kondisi tersebut juga terlihat pada data harga harian cabai rawit merah Kabupaten Sukabumi selama periode Januari 2025 hingga April 2026, di mana harga komoditas tersebut berubah dari sekitar Rp28.000/kg hingga mencapai Rp100.000/kg \parencite{data-pemerintah}. Fluktuasi harga tersebut menunjukkan bahwa komoditas pangan memiliki karakteristik yang dinamis sehingga diperlukan sistem yang mampu memantau sekaligus memprediksi pergerakan harga berdasarkan data historis.

Fluktuasi harga komoditas pangan yang tinggi menunjukkan bahwa informasi harga aktual saja belum cukup untuk mendukung proses pengambilan keputusan. Prediksi harga komoditas pertanian menjadi masukan penting bagi berbagai pemangku kepentingan, termasuk pemerintah, dalam merencanakan kebijakan, menjaga stabilitas pasokan, serta mendukung kesejahteraan masyarakat. Oleh karena itu, diperlukan suatu sistem yang mampu melakukan prediksi harga secara akurat berdasarkan data historis dan faktor-faktor yang memengaruhinya sehingga dapat menghasilkan informasi yang lebih bernilai dibandingkan sekadar penyajian data harga saat ini \parencite{bhardwaj2023innovativedeeplearningbased}.

Berdasarkan permasalahan tersebut, diperlukan suatu sistem yang dapat mendukung proses pemantauan harga komoditas pangan secara berkelanjutan melalui pemanfaatan data historis. Sistem tersebut diharapkan mampu mengintegrasikan proses pengumpulan data harga pasar, harga petani, serta data ketersediaan dan kebutuhan harian dari instansi terkait sebagai masukan bagi layanan prediksi harga. Hasil prediksi tersebut kemudian disajikan kepada Tim Pengendalian Inflasi Daerah (TPID) sebagai informasi pendukung dalam kegiatan pemantauan harga dan pengambilan keputusan. Oleh karena itu, dikembangkan Sistem Informasi Pemantauan Harga Komoditas (SIPANTAU) berbasis web sebagai media yang mengintegrasikan proses pengumpulan data dan penyajian hasil prediksi harga komoditas pangan di Kabupaten Sukabumi.

% Lorem ipsum dolor sit amet, consectetur adipiscing elit. Quisque
% tincidunt tellus sit amet tincidunt lobortis. Nulla id tempus
% enim. Curabitur sed enim quam. Duis suscipit in nibh nec
% gravida. Curabitur placerat ullamcorper dolor a rhoncus. Duis ac
% suscipit enim. Ut elementum tortor ligula, fringilla viverra diam
% varius eget. Vivamus quis iaculis lorem. Vivamus nec dignissim libero,
% eu congue sem. Nam id lobortis elit \parencite{foster}.

% Phasellus porttitor purus \textcite{warn} vitae finibus laoreet. Class
% aptent taciti sociosqu ad litora torquent per conubia nostra, per
% inceptos himenaeos. Lorem ipsum dolor sit amet, consectetur adipiscing
% elit. Nunc tempus sed purus sit amet porta. Sed volutpat est ac arcu
% porttitor, non pharetra tortor pharetra. Phasellus vitae dictum
% elit. Curabitur elementum consequat elementum. Integer in vulputate
% metus. Aliquam erat volutpat. Phasellus vitae convallis orci. Etiam ac
% metus tristique, viverra metus a, mattis lectus. Aenean vel vestibulum
% augue.

\section{Rumusan Masalah}

% Phasellus porttitor purus vitae finibus laoreet. Class aptent taciti
% sociosqu ad litora torquent per conubia nostra, per inceptos
% himenaeos. Lorem ipsum dolor sit amet, consectetur adipiscing
% elit. Nunc tempus sed purus sit amet porta. Sed volutpat est ac arcu
% porttitor, non pharetra tortor pharetra. Phasellus vitae dictum
% elit. Curabitur elementum consequat elementum. Integer in vulputate
% metus. Aliquam erat volutpat. Phasellus vitae convallis orci. Etiam ac
% metus tristique, viverra metus a, mattis lectus. Aenean vel vestibulum
% augue.

\begin{enumerate}
\item Bagaimana membangun website SIPANTAU yang mendukung proses pengumpulan data harga komoditas pangan dari instansi terkait sebagai dasar prediksi harga?
\item Bagaimana menyediakan hasil prediksi harga komoditas pangan untuk tiga hari ke depan melalui website SIPANTAU sebagai pendukung kegiatan pemantauan harga oleh TPID?
\end{enumerate}


\section{Tujuan}

\begin{enumerate}
\item Membangun website SIPANTAU yang mendukung proses pengumpulan data harga komoditas pangan dari instansi terkait sebagai dasar prediksi harga.
\item Menyediakan hasil prediksi harga komoditas pangan untuk tiga hari ke depan melalui website SIPANTAU sebagai pendukung kegiatan pemantauan harga oleh TPID.
\end{enumerate}


\section{Manfaat}

\begin{enumerate}
\item Pemerintah Kabupaten Sukabumi (Dinas terkait) \par
Membantu Dinas Pertanian, Dinas Perdagangan, dan Dinas Ketahanan Pangan dalam melakukan pengelolaan serta pembaruan data harga komoditas dan data pendukung secara terpusat sehingga dapat digunakan sebagai dasar proses prediksi harga komoditas pangan.
\item Tim Pengendali Inflasi Daerah \par
Memberikan informasi hasil prediksi harga komoditas pangan sebagai bahan pendukung dalam proses pemantauan harga serta pengambilan keputusan terkait pengendalian inflasi dan stabilisasi harga.
\item Penulis \par
Menambah pengalaman dalam pengembangan layanan backend berbasis web menggunakan framework Laravel, pengelolaan basis data MySQL, integrasi layanan prediksi, serta proses deployment aplikasi pada lingkungan server.
\end{enumerate}

% Curabitur elementum consequat elementum. Integer in vulputate
% metus. Aliquam erat volutpat. Phasellus vitae convallis orci. Etiam ac
% metus tristique, viverra metus a, mattis lectus. Aenean vel vestibulum
% augue.


\section{Batasan Masalah}

\begin{enumerate}
\item Sistem yang dibahas adalah Sistem Informasi Pemantauan Harga Komoditas (SIPANTAU) pada lingkup Kabupaten Sukabumi.
\item Data yang digunakan merupakan data harga pasar, harga petani, data ketersediaan dan kebutuhan harian, dan data panen komoditas pangan yang diperoleh dari Pemerintah Kabupaten Sukabumi sebagai data internal.
\item Komoditas yang dianalisis dalam sistem ini dibatasi pada beberapa komoditas strategis, yaitu cabai rawit merah, cabai merah besar, bawang merah, dan beras.
\item Implementasi sistem difokuskan pada wilayah Pasar Cisaat sebagai studi kasus selama periode pelaksanaan Praktik Kerja Lapang untuk memastikan sistem berjalan secara fungsional.
\item Sistem berfokus pada proses pengumpulan data, integrasi layanan prediksi harga, dan penyajian hasil prediksi, serta tidak mencakup pengembangan model prediksi maupun implementasi kebijakan pemerintah secara langsung.
\item Sistem menggunakan pembagian hak akses (\emph{role-based access control}), yaitu Dinas (Pertanian, Perdagangan, dan Ketahanan Pangan) sebagai pengelola dan input data, serta TPID sebagai pengguna fitur analisis dan prediksi.
\item Integrasi data eksternal melalui API (seperti SILINDA) masih dalam tahap pengembangan dan belum diimplementasikan pada sistem saat ini.
\item Model prediksi harga merupakan layanan yang telah tersedia dan tidak menjadi bagian dari pengembangan pada kegiatan Praktik Kerja Lapang ini. Ruang lingkup pengembangan difokuskan pada integrasi layanan prediksi ke dalam website SIPANTAU melalui backend.
\end{enumerate}


% \section{Sistematika Pembahasan}
% \noindent
% \textbf{BAB I : PENDAHULUAN}

% Pada bab ini dijelaskan latar belakang, rumusan masalah, batasan,
% tujuan, manfaat,  dan sistematika penulisan.\\

% \noindent
% \textbf{BAB II : TINJAUAN PUSTAKA DAN LANDASAN TEORI}

% Pada bab ini dijelaskan teori-teori dan penelitian terdahulu yang
% digunakan sebagai acuan dan dasar dalam penelitian.\\

% \noindent
% \textbf{BAB III : METODOLOGI PENELITIAN}

% Pada bab ini dijelaskan metode yang digunakan dalam penelitian
% meliputi langkah kerja, pertanyaan penilitian, alat dan bahan, serta
% tahapan dan alur penelitian.\\

% \noindent
% \textbf{BAB IV : HASIL DAN PEMBAHASAN}

% Pada bab ini dijelaskan hasil penelitian dan pembahasannya.\\

% \noindent
% \textbf{BAB V : KESIMPULAN DAN SARAN}

% Pada bab ini ditulis kesimpulan akhir dari penelitian dan saran untuk
% pengembangan penelitian selanjutnya.\\
